\documentclass[11pt]{amsart}

\usepackage[T1]{fontenc}
\usepackage{lmodern}
\usepackage{microtype}
\usepackage{amsmath,amssymb,amsthm}
\usepackage[hidelinks]{hyperref}

\hypersetup{
  pdftitle={A Counterexample at n=6 to Theorem 1 and Conjecture 36(3) of Kundu and Velmurugan}
}

\newtheorem{theorem}{Theorem}
\newtheorem{corollary}[theorem]{Corollary}
\theoremstyle{remark}
\newtheorem{remark}[theorem]{Remark}

\newcommand{\Irr}{\operatorname{Irr}}
\newcommand{\ccn}{\operatorname{ccn}}
\newcommand{\one}{\mathbf{1}}

\title[A counterexample at $n=6$]
{A Counterexample at $n=6$ to Theorem~1 and Conjecture~36(3)
of Kundu and Velmurugan}

\date{}

\subjclass[2020]{20C30, 05E10}
\keywords{
covering number, symmetric group, irreducible character,
Kronecker product
}

\begin{document}

\begin{abstract}
We give a short calculation showing that
\[
  \ccn(\chi_{(3,3)};S_6)
  =
  \ccn(\chi_{(2,2,2)};S_6)
  =5.
\]
Thus the upper inequality in Conjecture~36(3) of Kundu and
Velmurugan is false at \(n=6\); the same example also contradicts the
bound in their Theorem~1, which excludes neither \((3,3)\) nor
\((2,2,2)\). The value itself is implicit in existing formulas
for tensor powers of the irreducible representations of \(S_6\); what
is new here is the resulting contradiction with these two statements of
Kundu and Velmurugan.
\end{abstract}

\maketitle

\section{The counterexample}

For a character \(\theta\) of a finite group \(G\), let \(c(\theta)\)
denote its set of irreducible constituents. When it exists, define
\(\ccn(\theta;G)\) to be the least positive integer \(k\) such that
\[
  c(\theta^k)=\Irr(G),
\]
where \(\theta^k\) is the pointwise product, equivalently the character
of the \(k\)-fold tensor power. Sun, Zhang, and Zhu computed power
formulas for all irreducible representations of \(S_6\); their formulas
also imply the value below \cite{SZZ}. We give a short independent
class-function proof and record its consequence for \cite{KV}.

\begin{theorem}\label{thm:main}
For the irreducible characters of \(S_6\) indexed by \((3,3)\) and
\((2,2,2)\),
\[
  \ccn(\chi_{(3,3)};S_6)
  =
  \ccn(\chi_{(2,2,2)};S_6)
  =5.
\]
\end{theorem}

\begin{proof}
Set \(\chi=\chi_{(3,3)}\), and let \(\pi_r\) be the permutation
character of \(S_6\) acting on its \(r\)-subsets. Young's rule gives
\[
  \pi_3
  =
  \chi_{(6)}+\chi_{(5,1)}+\chi_{(4,2)}+\chi_{(3,3)},
  \qquad
  \pi_2
  =
  \chi_{(6)}+\chi_{(5,1)}+\chi_{(4,2)},
\]
so \(\chi=\pi_3-\pi_2\). If \(C_\mu\) is the conjugacy class of cycle
type \(\mu=(\mu_1,\ldots,\mu_t)\), then
\[
  \pi_r(\mu)
  =
  [x^r]\prod_{i=1}^t(1+x^{\mu_i}).
\]
Writing \(\varepsilon\) for the sign character, we obtain
\[
\begin{array}{c|rrrrrrrrrrr}
\mu
 &1^6&21^4&2^21^2&2^3&31^3&321&3^2&41^2&42&51&6\\ \hline
|C_\mu|
 &1&15&45&15&40&120&40&90&90&144&120\\
\chi(\mu)
 &5&1&1&-3&-1&1&2&-1&-1&0&0\\
\varepsilon(\mu)
 &1&-1&1&-1&1&-1&1&-1&1&1&-1
\end{array}
\]
The usual inner product of class functions therefore gives
\[
  \langle\chi,\one\rangle
  =
  \langle\chi^2,\varepsilon\rangle
  =
  \langle\chi^3,\one\rangle
  =
  \langle\chi^4,\varepsilon\rangle
  =0.
\]
For example,
\[
\begin{aligned}
720\langle\chi^4,\varepsilon\rangle
={}&625-15+45-1215+40-120\\
   &\quad+640-90+90
 =0.
\end{aligned}
\]
Hence none of \(\chi,\chi^2,\chi^3,\chi^4\) contains every irreducible
character of \(S_6\), and thus
\[
  \ccn(\chi;S_6)\geq5.
\]

Miller proved that, for \(n>4\), every nonlinear
\(\psi\in\Irr(S_n)\) satisfies
\[
  c(\psi^{n-1})=\Irr(S_n)
\]
\cite[Theorem~1.1(i)]{Miller}. Since \(\chi\) is nonlinear and
\(n-1=5\) for \(n=6\), this gives
\(c(\chi^5)=\Irr(S_6)\), and therefore
\[
  \ccn(\chi_{(3,3)};S_6)=5.
\]

Finally, \((3,3)'=(2,2,2)\), so
\[
  \chi_{(2,2,2)}=\chi_{(3,3)}\varepsilon.
\]
For every \(k\),
\[
  (\chi\varepsilon)^k=\chi^k\varepsilon^k.
\]
Tensoring by the linear character \(\varepsilon^k\) permutes
\(\Irr(S_6)\). Thus the two characters have the same covering number.
\end{proof}

\begin{corollary}\label{cor:consequences}
The upper inequality in Conjecture~36(3) of Kundu and Velmurugan
\cite{KV} is false. The bound in their Theorem~1 is also false as
stated.
\end{corollary}

\begin{proof}
Conjecture~36(3) asserts, for even \(n\geq5\),
\[
  \ccn\!\left(\chi_{(n/2,n/2)};S_n\right)
  \leq
  \lceil\log_2 n\rceil+1.
\]
At \(n=6\), the right-hand side is \(4\), whereas
Theorem~\ref{thm:main} gives
\[
  \ccn(\chi_{(3,3)};S_6)=5.
\]

Their Theorem~1 asserts, for \(n\geq5\) and
\[
  \lambda
  \notin
  \{(n),(1^n),(n-1,1),(2,1^{n-2})\},
\]
that
\[
  \ccn(\chi_\lambda;S_n)
  \leq
  \left\lceil\frac{2(n-1)}{3}\right\rceil.
\]
At \(n=6\), this bound is \(4\), and neither \((3,3)\) nor
\((2,2,2)\) is excluded. Theorem~\ref{thm:main} therefore contradicts
the stated bound.
\end{proof}

\begin{remark}
The lower inequality in Conjecture~36(3) remains valid at \(n=6\).
The counterexample does not address the corresponding upper bound
restricted to even \(n\geq8\). In the proof of their Theorem~1, Kundu
and Velmurugan state that the cases \(n\leq11\) were verified directly
using SageMath; the examples above show that this finite verification
cannot be correct as stated.
\end{remark}

\begin{thebibliography}{9}

\bibitem{KV}
R.~Kundu and S.~Velmurugan,
\emph{Covering numbers of some irreducible characters of the symmetric
group},
Electron. J. Combin. \textbf{32} (2025), no.~2,
Paper No.~P2.56, 39 pp.
\href{https://doi.org/10.37236/13289}{doi:10.37236/13289}.

\bibitem{Miller}
A.~R. Miller,
\emph{Covering numbers for characters of symmetric groups},
Ann. Sc. Norm. Super. Pisa Cl. Sci. (5)
\textbf{26} (2025), no.~1, 517--524.
\href{https://doi.org/10.2422/2036-2145.202302_003}
{doi:10.2422/2036-2145.202302\_003}.

\bibitem{SZZ}
J.-C. Sun, C. Zhang, and H. Zhu,
\emph{Kronecker coefficients and Harrison centres of the
representation ring of the symmetric group},
\href{https://arxiv.org/abs/2407.18152}
{arXiv:2407.18152}, 2024.

\end{thebibliography}

\end{document}